\begin{trapbox}
	\begin{description}[style=nextline, leftmargin=2.8em, labelsep=0.5em, itemsep=0.35em]
		\item[\textbf{\textcolor{CTrap}{易错点一（忽略轴位）：}}] 不判断对称轴是否在区间内，直接认为最值就在顶点 $f(x_0)$ 处取得。
		\item[\textbf{\textcolor{CTrap}{正确理解（先判位置）：}}] 必须判断 $x_0$ 是否属于 $[m,n]$：若属于，则顶点参与最值；若不属于，则顶点不在区间内，最值只能由端点函数值决定。
		\item[\textbf{\textcolor{CTrap}{错因分析（刻板记忆）：}}] 学生常误记“二次函数最值就是顶点值”，忽视区间限制，导致分类遗漏。
	\end{description}
	\medskip
	\begin{description}[style=nextline, leftmargin=2.8em, labelsep=0.5em, itemsep=0.35em]
		\item[\textbf{\textcolor{CTrap}{易错点二（开口颠倒）：}}] 混淆开口方向，比如 $a>0$ 却将顶点当做最大值，或 $a<0$ 将顶点当做最小值。
		\item[\textbf{\textcolor{CTrap}{正确理解（记开口本质）：}}] $a>0$ 开口向上，顶点是最小值点，最大值在离顶点较远的端点；$a<0$ 开口向下，顶点是最大值点，最小值在离顶点较远的端点。
		\item[\textbf{\textcolor{CTrap}{错因分析（符号混淆）：}}] 对二次项系数 $a$ 的符号意义不熟练，导致最值属性错配。
	\end{description}
	\medskip
	\begin{description}[style=nextline, leftmargin=2.8em, labelsep=0.5em, itemsep=0.35em]
		\item[\textbf{\textcolor{CTrap}{易错点三（端点比较机械）：}}] 在比较 $f(m)$ 和 $f(n)$ 时，单纯地认为“左边小右边大”或“右大左小”，未考虑对称轴位置影响函数值大小。
		\item[\textbf{\textcolor{CTrap}{正确理解（直接计算比较）：}}] 无论开口方向，比较最值时直接计算 $f(m), f(n)$ 并取较大/小即可，这是最稳妥的方法。
		\item[\textbf{\textcolor{CTrap}{错因分析（过度推理）：}}] 试图用单调性定位，但未正确判断单调区间，易出现反转错误；直接计算端点函数值可以避免此问题。
	\end{description}
\end{trapbox}
